\documentclass[12pt, letterpaper]{article}

%% ─── PACKAGES ───────────────────────────────────────────────────────────────
\usepackage[top=1in, bottom=1in, left=1in, right=1in]{geometry}
\usepackage{setspace}          % \doublespacing
\usepackage{titlesec}          % numbered section headings
\usepackage{fancyhdr}          % running header / footer
\usepackage{microtype}         % microtypographic refinements
\usepackage{enumitem}          % fine-grained list control
\usepackage{booktabs}          % professional table rules
\usepackage{tabularx}          % auto-width table columns
\usepackage{array}             % extended column types
\usepackage{listings}          % verbatim code blocks (search strings)
\usepackage[hidelinks]{hyperref}

%% ─── PARAGRAPH & SPACING ────────────────────────────────────────────────────
\setlength{\parindent}{0.5in}
\setlength{\parskip}{0pt}

%% ─── SECTION HEADING FORMAT (numbered, bold) ────────────────────────────────
\titleformat{\section}{\normalfont\large\bfseries}{\thesection.}{0.5em}{}
\titlespacing*{\section}{0pt}{18pt}{6pt}

%% ─── LIST DEFAULTS ──────────────────────────────────────────────────────────
\setlist{nosep, itemsep=4pt, parsep=0pt, topsep=6pt}

%% ─── LISTINGS SETUP for search strings ─────────────────────────────────────
%% Frame around the search block, monospace font, line-wrapping enabled.
\lstset{
  basicstyle  = \small\ttfamily,
  breaklines  = true,
  breakatwhitespace = false,
  frame       = single,
  framesep    = 6pt,
  xleftmargin = 0.4in,
  xrightmargin= 0.4in,
  aboveskip   = 8pt,
  belowskip   = 8pt,
  lineskip    = 2pt,
  keepspaces  = true
}

%% ─── HEADER / FOOTER ────────────────────────────────────────────────────────
\pagestyle{fancy}
\fancyhf{}
\lhead{\small Kamrul, Md Kamruzzaman}
\rhead{\small EDCI 59100-016 $|$ Summer 2026}
\cfoot{\thepage}
\renewcommand{\headrulewidth}{0.4pt}
\setlength{\headheight}{14pt}

%% ════════════════════════════════════════════════════════════════════════════
\begin{document}
%% ════════════════════════════════════════════════════════════════════════════

%% ─── TITLE BLOCK ────────────────────────────────────────────────────────────
\begin{center}
  {\Large\bfseries Initial Literature Review Plan}\\[0.75em]
  {\normalsize EDCI 59100-016 $|$ Summer 2026}\\[0.25em]
  {\normalsize Md Kamruzzaman Kamrul}\\[0.25em]
  {\normalsize May 31, 2026}
\end{center}

\bigskip
\doublespacing

%% ─────────────────────────────────────────────────────────────────────────────
\section{Topic}
%% ─────────────────────────────────────────────────────────────────────────────

This review examines how Vision-Language Models (VLMs) and multimodal
architectures have been applied to spatiotemporal reasoning for safety
monitoring and activity tracking on active construction sites.

%% ─────────────────────────────────────────────────────────────────────────────
\section{Research Questions and Purpose}
%% ─────────────────────────────────────────────────────────────────────────────

\noindent\textbf{Purpose:} To systematically map the current state of
VLM research applied to spatiotemporal reasoning in construction---
identifying prevalent architectures, data fusion strategies, deployment
barriers, and evaluation frameworks---and to locate the gaps where future
primary research is most needed.

\noindent\textbf{Research Questions:}

\begin{itemize}[leftmargin=0.75in, label={}]
  \item \textbf{RQ1:} What are the prevalent VLM architectures currently
        utilized for spatiotemporal analysis in construction environments?
  \item \textbf{RQ2:} How is multimodal data---visual and textual---fused
        and processed to evaluate safety hazards on active construction sites?
  \item \textbf{RQ3:} What are the primary methodological bottlenecks in
        deploying these models for real-time progress tracking and safety
        monitoring?
  \item \textbf{RQ4:} How is the performance of VLMs currently evaluated and
        benchmarked for spatiotemporal tasks in construction contexts?
\end{itemize}

%% ─────────────────────────────────────────────────────────────────────────────
\section{Selected Review Methodology}
%% ─────────────────────────────────────────────────────────────────────────────

\noindent\textbf{Methodology:} Scoping review, structured according to
PRISMA-ScR (Preferred Reporting Items for Systematic Reviews and Meta-Analyses
extension for Scoping Reviews) guidelines.
This methodology maps the breadth and nature of available evidence---
architectures, datasets, and evaluation practices---without requiring
methodological homogeneity across included studies. It is appropriate for an
emerging, multidisciplinary field in which evaluation metrics and study
designs vary substantially across the literature, making statistical
meta-analysis premature.

%% ─────────────────────────────────────────────────────────────────────────────
\section{Search Strategy}
%% ─────────────────────────────────────────────────────────────────────────────

\noindent\textbf{Keyword development.} Search keywords were identified by
analyzing terminology used in recent studies that apply artificial intelligence
to construction management and safety. The search query was then structured
using a tripartite Boolean logic framework designed to intersect three distinct
conceptual domains: (1)~the specific AI architecture (multimodal vision-language
systems); (2)~the analytical task (spatiotemporal and dynamic reasoning); and
(3)~the application domain (construction management and safety).

\noindent\textbf{Search strings applied across both databases:}

\begin{singlespace}
\begin{lstlisting}[caption={},nolol]
String 1 - AI Architecture:
("Vision-Language Model*" OR "VLM" OR "Vision-Text"
OR "Image-Text" OR "Visual Question Answering" OR "VQA"
OR "Image Captioning" OR "CLIP" OR "Large Multimodal Model*"
OR "LMM" OR "Foundation Model*" OR "Generative AI"
OR "Video-Language" OR "Video-Text")

AND

String 2 - Analytical Task:
("Spatiotemporal" OR "Spatio-temporal" OR "Video"
OR "Temporal" OR "Dynamic" OR "Time-series" OR "4D"
OR "Motion" OR "Tracking" OR "Sequenc*"
OR "Action recognition" OR "Activity recognition")

AND

String 3 - Application Domain:
("Construction Management" OR "Civil Engineering" OR "AEC"
OR "Infrastructure" OR "Built Environment"
OR "Construction Site" OR "Safety Monitoring"
OR "Structural Health Monitoring" OR "Digital Twin")
\end{lstlisting}
\end{singlespace}

\noindent\textbf{Refinement plan.} After running the initial string in each
database and recording hit counts, result relevance will be evaluated. If the
hit set is too broad---for example, dominated by general computer vision papers
with no construction context---subject-area filters (engineering, civil
engineering, computer science) will be applied within the database interface.
If the hit count is too low, additional synonyms will be introduced within the
relevant cluster (e.g., \textbf{jobsite}, \textbf{worksite} alongside
\textbf{Construction Site}).
The search is also supplemented by a \textbf{backward snowball phase}---
examining reference lists of included full-text papers---to capture foundational
computer vision and early multimodal studies that predate the core VLM search
window but remain essential comparative baselines.

\noindent\textbf{Anticipated challenges.} Because this topic spans construction
management, computer vision, and natural language processing, relevant papers
may be distributed across journals that do not fully overlap in any single
database. Terminology is also inconsistent across disciplines: the same class
of model may appear as a \textbf{vision--language model}, \textbf{multimodal
transformer}, or \textbf{large multimodal model} in different papers. The broad
synonym coverage in String~1 is designed specifically to address this risk.
All synonym expansions and any deviations from the original strings will be
documented throughout the process.

%% ─────────────────────────────────────────────────────────────────────────────
\section{Databases and Sources}
%% ─────────────────────────────────────────────────────────────────────────────

The primary databases are \textbf{Web of Science} (Core Collection) and
\textbf{Scopus}. These two platforms were selected because together they
provide broad, multidisciplinary coverage of high-impact peer-reviewed journals
spanning civil engineering, construction management, computer science, and
applied AI---the four disciplines most directly relevant to this review.

Web of Science is particularly strong in journals indexed in the Science
Citation Index Expanded (SCIE), which includes leading safety science,
automation, and engineering venues. Scopus complements this with wider coverage
of applied engineering and construction management journals, and is frequently
cited in scoping review methodology as a preferred companion database to Web
of Science. Running the identical search string across both platforms in
parallel reduces the risk that relevant literature indexed in only one database
will be missed---a recognized best practice in evidence synthesis.

Google Scholar will not serve as a primary database because its results are not
consistently reproducible across sessions, which conflicts with the PRISMA-ScR
transparency requirements. Preprint servers such as arXiv are also excluded,
as the review is restricted to high-impact peer-reviewed publications.

%% ─────────────────────────────────────────────────────────────────────────────
\section{Inclusion and Exclusion Criteria}
%% ─────────────────────────────────────────────────────────────────────────────

The eligibility criteria below are organized along four dimensions that directly
correspond to the scope implied by the four research questions. These criteria
are applied sequentially at both the title-and-abstract screening stage and the
full-text review stage.

\vspace{8pt}
%% ── Inclusion / Exclusion Table ─────────────────────────────────────────────
\begin{singlespace}
\small
\renewcommand{\arraystretch}{1.55}
\begin{tabularx}{\textwidth}{>{\bfseries}p{2.3cm} X X}
\toprule
\textbf{Dimension} & \textbf{Inclusion Criteria} & \textbf{Exclusion Criteria} \\
\midrule
Publication Date
  & 2019--2026 for core VLM literature; 2010--2026 for foundational
    computer vision baselines.
  & Pre-2010; post-March 2026. \\
\midrule
Source Type
  & High-impact, peer-reviewed journal articles.
  & Editorials, book chapters, theses, non-English documents,
    conference proceedings, and industry white papers. \\
\midrule
Technology
  & Multimodal architecture integrating visual/geometric data with
    textual/semantic data (e.g., VLMs, VQA, Image Captioning,
    CLIP-based models, Large Multimodal Models).
  & Purely unimodal computer vision (e.g., standalone YOLO, CNN
    bounding-box detectors) or purely text-based NLP, unless serving
    as an explicit comparative baseline within a hybrid multimodal
    framework. \\
\midrule
Application Domain
  & Active construction phase: Safety Management (hazard identification,
    PPE compliance) and Progress Tracking (productivity analysis,
    schedule monitoring, reporting).
  & Post-construction or facility management contexts, structural health
    monitoring outside the active construction phase, and
    non-engineering domains. \\
\bottomrule
\end{tabularx}
\end{singlespace}

\vspace{8pt}
\noindent\textbf{Screening procedure.} Screening will follow a two-stage
PRISMA-ScR process. Stage~1 screens titles and abstracts to remove clearly
out-of-scope records. Stage~2 applies the full eligibility criteria to
remaining full texts, with one standardized exclusion reason recorded per
rejected paper (e.g., \textbf{wrong domain}, \textbf{non-peer-reviewed},
\textbf{no temporal component}, \textbf{purely unimodal}). Record counts at
each stage will be reported in a PRISMA-ScR flow diagram included in the
final submission.

%% ─────────────────────────────────────────────────────────────────────────────
\section{Reflection on the Search Process}
%% ─────────────────────────────────────────────────────────────────────────────

Three developments have meaningfully shaped this plan since the initial
proposal. First, in response to instructor feedback from Week~1, the original
broad framing---VLMs in built environments generally---was narrowed to focus
specifically on safety monitoring and activity tracking during the active
construction phase; this change made the search terms, database choices, and
eligibility criteria substantially more coherent and focused. Second, a single
broad research question was restructured into four operationally distinct
questions (RQ1--RQ4) that define clear, non-overlapping mapping targets:
the Technology inclusion criterion now directly operationalizes the scope
implied by RQ1 and RQ2, while the Application Domain criterion reflects the
construction-specific focus required by all four questions. Third, the
adoption of PRISMA-ScR as the reporting standard formalized what were initially
informal methodological decisions, and the inclusion of a backward snowball
phase emerged from recognizing that foundational computer vision baselines
(2010 onward) often predate the keyword-indexed VLM literature and would
otherwise be missed by the primary database search alone.

%% ════════════════════════════════════════════════════════════════════════════
\end{document}
%% ════════════════════════════════════════════════════════════════════════════